\documentclass[ ../main.tex]{subfiles}
\providecommand{\mainx}{..}
\begin{document}
\section{$k$-perfect hash functions}
An $k$-perfect hash function permits up to $k$ collisions and is therefore a \emph{generalization} of the hashing family, i.e., a \emph{perfect hash function} is a special case where $k=0$.

\begin{algorithm}
	\caption{$k$-perfect hash function generator}
	\label{alg:kph}
	\SetKwProg{func}{function}{}{}
	\KwIn
	{
		$\Set{S}$ is a set, $r$ is a load factor, and $k$ is the maximum number of collisions allowed.
	}
	\KwOut
	{t
		A $\ell$-perfect hash function of $\Set{S}$ with a load factor $r$ and $\ell \leq k$.
	}
	\func{\makehash{$\Set{S}$, $r$, $k$}}
	{
		$m \gets \Card{\Set{S}}$\;
		$N \gets \left\lceil \frac{m}{r} \right\rceil$\;
		\For{$n \gets 1$ \KwTo $\infty$}
		{
			\tcp{We give a name to the hash function that is coded by
				$n'$. Note that $\catop$ is the concatenation operator.}
			$\hat{\hash}(x) \equiv \hash(x \catop n') \mod N$\;
			$\Set{H} \gets \EmptySet$\;
			$\ell \gets 0$\;
			\For{$x \in \Set{S}$}
			{
				$h \gets \hat{\hash}(x)$\;
				\If{$h \in \Set{H}$}
				{
					$\ell \gets \ell + 1$\;
				}
				$\Set{H} \gets \Set{H} \cup \{ h \}$\;
			}
			\If{$\ell \leq k$} 
			{
				\tcp{This tuple is sufficient to code the \emph{$k$-perfect hash function} for $\Set{S}$ that has the minimum bit length.}
				\Return $(n',N)$\;
			}                
		}        
	}
\end{algorithm}


The probability that $t$ collisions occur for a particular bit string in \cref{alg:kph} is given by the following theorem.
\begin{theorem}
	The probability that a bit string $b$ results in a $t$-perfect hash function $\hash_{\Set{S}}$ in \cref{alg:kph} is given by
	\begin{equation}
	\label{eq:p_m_N}
	p(m,N) = \frac{N!}{N^m (N-m)!}\,,
	\end{equation}
	where $m = \Card{\Set{S}}$ and $N$ is the maximum hash of $\hash_{\Set{S}}$.
\end{theorem}
\begin{proof}
	Suppose we have a set $\Set{S} = \{x_1,\ldots,x_m\}$, $t \geq m$. The probability that $x_1,\ldots,x_t$ collide and the rest $x_{t+1},\ldots,x_m$ do not collide is given by
	\begin{equation}
	\frac{1}{N^t} \prod_{i=1}^{m-t} \frac{N - i}{N} = \frac{(N-1)!}{N^m(N-m+t-1)!}\,.
	\end{equation}
	We are actually interested in \emph{any} $t$ elements $\{x_{j_1},\ldots,x_{j_t}\}$ colliding, not just $x_1,\ldots,x_t$. There are $m \choose t$ ways to choose $t$ of elements in $\Set{S}$ as a collision. Thus, the probability that ...
	\begin{equation}
	p(m,t,N) = \frac{m!}{t!(m-t)!} \frac{(N-1)!}{N^m(N-m+t-1)!}\,.
	\end{equation}
	
	We are interested in the first case when $t$ collisions occur, which is a geometric distribution with probability of success $p(m,t,N)$ as given by the discrete random variable
	\begin{equation}
	\RV{Q} \sim \geodist\!\bigl(p\!\left(m,t,N\right)\bigr)\,.
	\end{equation}
	The expected number of trials for the geometric distribution is given by
	\begin{equation}
	\Expect{\RV{Q}} = \frac{1}{p(m,t,N)} = \frac{t!(m-t)!N^m(N-m+t-1)!}{m!(N-1)!}\,.
	\end{equation}
	
	TODO: Taylor approximation. Should work out well...but maybe not worth doing...i.e., this is close enough?
	
	By \cref{def:mapping}, the \nth trial uniquely maps to a bit string of length $m = \lfloor \log_2 n \rfloor$. Thus, the expected bit length is given approximately by
	\begin{align}
	\Expect{\log_2 \RV{Q}}
	&= \log_2\!\left(\frac{t!(m-t)!N^m(N-m+t-1)!}{m!(N-1)!}\right)\\
	&= \log_2 t! + \log_2(m-t)! + m \log_2\! N +
	\log_2\!\left(N-m+t-1\right)! - \log_2 m! - \log_2(N-1)! \; \si{bits}\,.
	\end{align}
	By Stirling's approximation, 
	\begin{equation}
	\log_2 n! \approx n \left(\log_2 n - \log_2 \mathrm{e}\right)\,.
	\end{equation}
	and so the expectation may be rewritten approximately as
	\begin{equation}
	\begin{split}
	\Expect{\RV{Q}} \approx \;
	&m \log_2 N + t \left(\log_2 t - \log_2 \mathrm{e}\right) +\\
	&(m - t) \left(\log_2 (m - t) - \log_2 \mathrm{e}\right) +\\
	&(N - m + t - 1) \left(\log_2 (N - m + t - 1) - \log_2 \mathrm{e}\right) -\\
	&m \left(\log_2 m - \log_2 \mathrm{e}\right) -\\
	&(N-1) \left(\log_2 (N-1) - \log_2 \mathrm{e}\right) \; \si{bits}\,.
	\end{split}
	\end{equation}
	
	\begin{equation}
	\begin{split}
	\Expect{\RV{Q}} \approx \;
	&m \log_2 N + t \log_2 t  + (m - t) \log_2 (m - t) +\\
	&(N - m + t - 1) \log_2 (N - m + t - 1) -\\
	&m \log_2 m - (N-1) \log_2 (N-1) - (t - m) \log_2 \mathrm{e} \; \si{bits}\,.
	\end{split}
	\end{equation}
	
	Dividing by $m$ results in
	\begin{equation}
	\begin{split}
	\Expect{\RV{Q}} \approx \;
	&\log_2 N + \frac{t}{m} \log_2 t  + \left(1 - \frac{t}{m}\right) \log_2 (m - t) +\\
	&(\frac{N}{m} - 1 + \frac{t}{m} - \frac{1}{m}) \log_2 (N - m + t - 1) -\\
	&\log_2 m - \frac{N-1}{m} \log_2 (N-1) - \left(\frac{t}{m} - 1\right) \log_2 \mathrm{e} \; \si{bits}\,.
	\end{split}
	\end{equation}
	
	Letting $\eta = \frac{t}{m}$ yields
	% n = t\m -> m n = t
	\begin{equation}
	\begin{split}
	\Expect{\RV{Q}} \approx \;
	&\log_2 N + \eta \log_2 m \eta  + \left(1 - \eta\right) \log_2 (m(1 - \eta)) +\\
	&(\frac{N}{m} - 1 + \eta - \frac{1}{m}) \log_2 (N - m(1 + \eta) - 1) -\\
	&\log_2 m - \frac{N-1}{m} \log_2 (N-1) - \left(\eta - 1\right) \log_2 \mathrm{e} \; \si{bits}\,.
	\end{split}
	\end{equation}
	
	Letting $r = \frac{m}{N}$ yields
	% r = m / N -> N = m / r
	\begin{equation}
	\begin{split}
	\Expect{\RV{Q}} \approx \;
	&\log_2\left(\frac{m}{r}\right) + \eta \log_2 m \eta  + \left(1 - \eta\right) \log_2 (m(1 - \eta)) +\\
	&\left(\frac{1}{r} - 1 + \eta - \frac{1}{m}\right) \log_2\left(\frac{m}{r} - m(1 + \eta) - 1\right) -\\
	&\log_2 m - \frac{\frac{m}{r}-1}{m} \log_2\left(\frac{m}{r}-1\right) - \left(\eta - 1\right) \log_2 \mathrm{e} \; \si{bits}\,.
	\end{split}
	\end{equation}
	
	Simplifying more yields
	\begin{equation}
	\begin{split}
	\Expect{\RV{Q}} \approx \;
	&- \log_2 r + \eta \log_2 \eta + \log_2 m + \log_2(1 - \eta) - \eta \log_2(1 - \eta) +\\
	&\left(\frac{1}{r} - 1 + \eta - \frac{1}{m}\right) \log_2\left(m\left(\frac{1}{r} - 1 - \eta\right) - 1\right) -\\
	&\left(\frac{1}{r}-\frac{1}{m}\right) \log_2\left(\frac{m}{r}-1\right) + \left(1 -\eta\right) \log_2 \mathrm{e} \; \si{bits}\,.
	\end{split}
	\end{equation}
	
	
	
	Simplifying more yields
	\begin{equation}
	\begin{split}
	\Expect{\RV{Q}} \approx \;
	&- \log_2 r + \eta \log_2 \eta + \log_2 m + \log_2(1 - \eta) - \eta \log_2(1 - \eta) +\\
	&\left(\frac{1}{r} - 1 + \eta - \frac{1}{m}\right) \log_2\left(m\left(\frac{m - m r - \eta m r - r}{m r}\right)\right) -\\
	&\left(\frac{1}{r}-\frac{1}{m}\right) \log_2\left(\frac{m}{r}-1\right) + \left(1 -\eta\right) \log_2 \mathrm{e} \; \si{bits}\,.
	\end{split}
	\end{equation}
	
	Here:
	\begin{equation}
	\begin{split}
	\Expect{\RV{Q}} \approx & \log_2 \mathrm{e} - \eta  \log_2 (1-\eta) + \eta \log_2 \eta - \\
	&\qquad \eta \log_2 \mathrm{e} + \log_2(1-\eta ) + \eta \log_2(\eta  m-1) -\\ &\qquad \frac{\log_2(\eta  m-1)}{m}+\frac{\log_2(m-1)}{m} + \log_2\left(\frac{m}{m-1}\right)
	\end{split}
	\end{equation}
\end{proof}
\end{document}